\chapter{Balanced evaluations}
\label{chap:balance}

\section{Preserving the surrounding stack}
% Introduce the source paper's definition: equal boundary stacks, with every
% intermediate stack extending the initial stack. Explain why boundary equality
% alone is insufficient for reasoning about evaluation under a continuation.

\section{Black-holing and update frames}
% Relate the big-step black-hole set to the machine stack by L = update_locs St.
% Show how an update frame extends this set and how other frames leave it alone.
% Distinguish an invariant derived from machine reachability from a side
% condition imposed directly by the balanced tree relation.

\section{The difficulty with linear traces}
% Use variable lookup and heap update to exhibit the required trace split:
% evaluation under an update frame, followed by the transition that removes it.
% Explain why endpoint-only and stack-indexed linear relations did not expose
% the subevaluations needed by the reverse induction.

\section{Balanced evaluation trees}
% Define the mutually inductive expression and matching tree relations.
% Show how their constructors record transitions and nested subevaluations.
% Explain why matching trees allow distinct initial and final argument stacks.
% Make clear that trees are not yet proved equivalent to all balanced sequences
% in the source paper's sense. The next chapter states exactly what is checked.
